\documentclass[12pt,a4paper]{article}
\usepackage[utf8]{inputenc}
\usepackage[vietnamese]{babel}
\usepackage[left=2cm,right=2cm,top=2cm,bottom=2cm]{geometry}
\usepackage{mathptmx}
\usepackage{enumitem}

\pagestyle{plain}
\linespread{1.15}

\begin{document}

% --- HEADER QUỐC HIỆU ---
\noindent
\begin{minipage}[t]{0.45\textwidth}
    \begin{center}
        \small
        CÔNG AN THÀNH PHỐ HÀ NỘI\\
        \textbf{\underline{PHÒNG KỸ THUẬT HÌNH SỰ}}\\[0.1cm]
        Số: 12/BCPY-BS\\
        \textit{V/v giám định pháp y bổ sung khung giờ tử vong}
    \end{center}
\end{minipage}
\hfill
\begin{minipage}[t]{0.52\textwidth}
    \begin{center}
        \small
        \textbf{CỘNG HÒA XÃ HỘI CHỦ NGHĨA VIỆT NAM}\\
        \textbf{\underline{Độc lập - Tự do - Hạnh phúc}}\\[0.1cm]
        \textit{Hà Nội, ngày 25 tháng 07 năm 2026}
    \end{center}
\end{minipage}

\vspace{0.6cm}

\begin{center}
    \textbf{\large BÁO CÁO GIÁM ĐỊNH PHÁP Y BỔ SUNG}\\
    \textit{(Phân tích nút thắt thời gian tử vong quyết định trong vụ án Khang)}
\end{center}

\vspace{0.3cm}

\section*{I. KẾT QUẢ PHÂN TÍCH HÓA SINH VÀ NỒNG ĐỘ MÁU}
\begin{enumerate}[label=\arabic*., leftmargin=1.5cm]
    \item \textbf{Khung giờ tử vong chính xác:} Dựa trên độ hoại tử tế bào cơ và hóa sinh dịch nhãn cầu, thời điểm nạn nhân Nguyễn Văn Khang bị đâm đứt động mạch cảnh tử vong xác định chính xác từ \textbf{21:00 đến 21:15} ngày 24/07/2026.
    \item \textbf{Đánh giá tổn thương gáy trước đó:} Vết va đập gáy xảy ra vào lúc khoảng 20:00 chỉ gây ngất bất tỉnh kéo dài khoảng 30–45 phút, không phải nguyên nhân tử vong.
    \item \textbf{KẾT LUẬN NÚT THẮT CHỨNG CỨ NGOẠI PHẠM:} Lúc 20:15 (khi Tùng tháo chạy trên camera), nạn nhân Khang vẫn còn sống trong trạng thái bất tỉnh. Hung thủ thực sự ra tay đâm chết nạn nhân vào lúc **21:00** (45 phút sau khi Tùng đã rời khỏi hiện trường).
\end{enumerate}

\section*{II. LỜI KHAI LỠ LỜI TRONG BIÊN BẢN THẨM VẤN BỔ SUNG TRẦN THỊ HÀ}
\noindent Khi bị triệu tập lần 2 đối chất về khung giờ tử vong 21:00, Trần Thị Hà biến đổi sắc mặt lúng túng và bộc lộ sơ hở chết người:\\
\textbf{ĐTV Lê Minh:} \textit{"Báo cáo pháp y khẳng định lúc 20:15 Khang chỉ mới bị ngất bất tỉnh và bị đâm đứt động mạch cảnh vào lúc 21:00. Chị khẳng định ở nhà cả tối không sang, vậy tại sao chị lại biết chính xác nạn nhân nằm gục ngay bên cạnh BỘ BÌNH TRÀ VỠ VỤN dưới sàn?"}\\
\textbf{Trần Thị Hà:} *(Ấp úng, biến đổi sắc mặt, lúng túng)* \textit{"Tôi... tôi... lúc lén ghé qua ngõ thấy nhà mở cửa, tôi chui vào định tìm anh ấy thì thấy anh Khang nằm bất tỉnh bên bộ bình trà vỡ... Tôi thề là tôi không đâm anh ấy!"}\\[0.1cm]
$\rightarrow$ \textbf{KẾT LUẬN LẠY ÔNG TÔI Ở BỤI NÀY:} Bộ bình trà chỉ bị Tùng làm vỡ lúc 20:00. Hà thừa nhận nhìn thấy bộ bình trà vỡ chứng tỏ Hà đã trực tiếp có mặt tại phòng khách sau 20:15 và ra tay sát hại nạn nhân lúc 21:00!

\vspace{0.8cm}

\noindent
\begin{minipage}[t]{0.45\textwidth}
    \small
    \textbf{\textit{Nơi nhận:}}\\
    - Trưởng Ban Chuyên án;\\
    - Viện KSMND TP. Hà Nội;\\
    - Lưu: VT, KTHS.
\end{minipage}
\hfill
\begin{minipage}[t]{0.5\textwidth}
    \begin{center}
        \small
        \textbf{GIÁM ĐỊNH VIÊN PHÁP Y}\\
        \textit{(Ký, ghi rõ họ tên)}\\[1.5cm]
        \textbf{Bác sĩ Chuyên khoa II Đỗ Văn Hùng}
    \end{center}
\end{minipage}

\end{document}
